\notepage{17}{From sampled functions to a continuum inner product}
Consider a function $f(x)$ on $a\le x\le b$, sampled at $n$ points
$x_1,\ldots,x_n$. Treat its samples as a coordinate column:
\[
 [\ket{f_n}]=\begin{bmatrix}f(x_1)\\f(x_2)\\\vdots\\f(x_n)\end{bmatrix}.
\]
\begin{center}
\begin{tikzpicture}[>=Stealth,x=1cm,y=1cm]
\draw[->] (0,0)--(6,0) node[right]{$x$};
\draw[->] (0,0)--(0,1.5) node[above]{$f(x)$};
\foreach \x/\lab in {1/x_1,2/x_2,4/x_i,5/x_n}
 {\draw (\x,-.06)--(\x,.06);\node[below] at (\x,-.08){$\lab$};}
\node at (3,.1){$\cdots$};
\node[above] at (1,0){$a$};\node[above] at (5,0){$b$};
\end{tikzpicture}
\end{center}
Introduce an orthonormal discrete coordinate basis $\ket{x_i}$ with
\[
 \braket{x_i}{x_j}=\delta_{ij},\qquad
 \sum_{i=1}^{n}\ket{x_i}\bra{x_i}=I.
\]
For example, the coordinate column of $\ket{x_\ell}$ has a 1 in
position $\ell$ and zeros elsewhere:
\[
 [\ket{x_\ell}]=
 \begin{bmatrix}\braket{x_1}{x_\ell}\\\vdots\\\braket{x_\ell}{x_\ell}\\\vdots\\\braket{x_n}{x_\ell}\end{bmatrix}
 =\begin{bmatrix}0\\\vdots\\1\\\vdots\\0\end{bmatrix}.
\]
The corresponding expansions and inner product are
\begin{align*}
 \ket{g_n}&=\sum_i\ket{x_i}\braket{x_i}{g_n}
           =\sum_i g(x_i)\ket{x_i},\\
 \bra{f_n}&=\sum_i f(x_i)^*\bra{x_i},\\
 \braket{f_n}{g_n}&=\sum_i f(x_i)^*g(x_i).
\end{align*}
As the sampling grows denser, the unweighted sum need not have a
finite limit. Introduce the Riemann-sum weight $\Delta_n=(b-a)/n$:
\[
 \braket{f_n}{g_n}_{\Delta}=\sum_i f(x_i)^*g(x_i)\Delta_n.
\]
For suitable functions the continuum limit is
\[
 \boxed{\braket{f}{g}=\lim_{n\to\infty}
 \sum_i f(x_i)^*g(x_i)\Delta_n=\int_a^b f(x)^*g(x)\dd x.}
\]
\editorial{The weighted inner product is a new convention; the
associated coordinate basis must be rescaled consistently in the
continuum limit. Note 18 gives its Dirac-delta normalization.}
